\chapter{Fundamental}


\section{YOLO (You Only Look Once)}

\subsection{Introduction}

The You Only Look Once (YOLO) framework, pioneered by Redmon et al. \cite{redmon2016you}, stands as a seminal contribution to the field of computer vision, specifically targeting the task of object detection within both static images and dynamic video sequences. A defining characteristic of YOLO is its ability to perform object detection in a remarkably fast and efficient manner through an end-to-end, single-pass methodology. This integrated approach streamlines the entire detection pipeline, leading to a compelling combination of high processing speed and competitive performance accuracy. Unlike earlier multi-stage detectors, YOLO processes the entire image at once, enabling it to reason globally about the image context and thus reducing the occurrence of background false positives.

The input to the YOLO model typically consists of an image or a frame extracted from a video stream, which is generally resized to a fixed, standardized dimension to ensure consistent processing. The output of the model is a comprehensive list of bounding boxes that precisely localize the detected objects within the input. For each identified bounding box, the model provides crucial information, including the coordinates of its center, its width and height, a confidence score reflecting the reliability of the detection, and the class label assigned to the detected object. This holistic output for each detected instance allows for a rich understanding of the objects present in the scene.

\subsection{Operational Principles}

YOLO employs a sophisticated approach based on mathematical principles and a deep convolutional neural network (CNN) architecture to achieve rapid and effective object detection in images. At the core of its operation, YOLO divides the entire input image into an $S \times S$ grid of equally sized cells. Each of these grid cells is tasked with the responsibility of predicting a fixed number, $B$, of bounding boxes and the associated object class probabilities for objects whose centers fall within that specific grid cell. This grid-based approach allows for simultaneous prediction across the entire image.

Each predicted bounding box is parameterized by a vector that encapsulates its spatial attributes and confidence:
$$
(b_x, b_y, b_w, b_h, C)
$$
where:
\begin{itemize}
    \item $b_x, b_y$: Represent the coordinates of the center of the bounding box relative to the bounds of the grid cell. These values are normalized to fall within the range of 0 to 1. This normalization ensures that the center coordinates are always within the predicting grid cell.
    \item $b_w, b_h$: Denote the width and height of the predicted bounding box, respectively. These dimensions are normalized by the overall width and height of the input image, ensuring scale invariance and allowing the model to predict boxes of varying sizes relative to the image.
    \item $C$: Represents the confidence score of the bounding box. This score reflects the model's certainty that an object is present within the bounding box and the accuracy of the predicted box. Mathematically, the confidence score is calculated as the product of the probability of an object being present within the bounding box ($P_{object}$) and the Intersection over Union (IoU) between the predicted bounding box and the ground truth bounding box ($IoU_{pred, truth}$):
    $$
    C = P_{object} \times IoU_{pred, truth}
    $$
    where the Intersection over Union (IoU) is defined as:
    $$
    IoU = \frac{\text{Area of Overlap}}{\text{Area of Union}}
    $$
    Here,
    \begin{itemize}
        \item Area of Overlap: Is the area of the intersection between the predicted bounding box and the ground truth bounding box. A larger overlap indicates a more accurate prediction.
        \item Area of Union: Is the total area encompassed by both the predicted and the ground truth bounding boxes, including the overlapping region. The IoU provides a normalized measure of the overlap, ranging from 0 (no overlap) to 1 (perfect overlap).
    \end{itemize}
\end{itemize}

In addition to predicting bounding boxes, each grid cell also predicts a conditional class probability vector, $P(c \mid object)$, where $c$ represents the possible object classes that the model is trained to detect. This vector provides the probability that the detected object, if present within the grid cell, belongs to a specific class among the predefined set of classes. This conditional probability is specific to the presence of an object within the grid cell.

The final class-specific confidence score for each bounding box is obtained by multiplying the individual bounding box confidence score, $C$, with the conditional class probability, $P(c \mid object)$:
$$
P(c) = P_{object} \times P(c \mid object)
$$
This score provides a comprehensive measure of the likelihood that a specific object of class $c$ is present within the bounding box and the model's overall confidence in this prediction.

Following the prediction phase, YOLO employs a crucial post-processing step known as Non-Maximum Suppression (NMS). The NMS algorithm is applied to filter out redundant and low-confidence bounding boxes, ensuring that each distinct object in the image is represented by a single, high-confidence bounding box. This process involves the following steps:
\begin{enumerate}
    \item For each class, all predicted bounding boxes with a confidence score below a certain threshold are discarded.
    \item The remaining bounding boxes are sorted in descending order based on their confidence scores.
    \item The bounding box with the highest confidence score is selected as a final detection.
    \item All other bounding boxes that have a significant overlap (above a predefined IoU threshold) with the selected box are suppressed (removed) to eliminate duplicates for the same object.
    \item Steps 3 and 4 are repeated until no more bounding boxes remain for that class.
\end{enumerate}
This NMS procedure is essential for reducing the number of false positive detections and refining the output to provide a more accurate and concise set of detection results, ensuring that each detected object has a unique and highly confident bounding box.


\section{DeepSORT (Deep Simple Online and Realtime Tracking)}

\subsection{Introduction}

DeepSORT (Deep Simple Online and Realtime Tracking), introduced by Wojke et al. \cite{wojke2017simple}, represents a significant evolution in the domain of Multi-Object Tracking (MOT) within video sequences. Building upon its predecessor, SORT (Simple Online and Realtime Tracking) \cite{bewley2016simple}, DeepSORT addresses inherent limitations, particularly in scenarios involving occlusions, changes in appearance, and long-term tracking of objects. The core innovation of DeepSORT lies in its seamless integration of deep learning-based appearance features, extracted through a dedicated deep neural network, into the tracking pipeline. This integration significantly enhances the system's robustness and accuracy in maintaining the identities of tracked objects across consecutive frames, even when faced with challenging visual variations.

The DeepSORT framework takes as input the detected bounding boxes of potential objects within each video frame, along with their associated confidence scores, typically generated by a state-of-the-art object detection model such as YOLO \cite{redmon2016you} or Faster R-CNN \cite{ren2015faster}. The primary output of DeepSORT is a set of trajectories, each uniquely identifying a tracked object throughout the video sequence. Each trajectory comprises the bounding box coordinates of the object in every frame where it is detected and a consistent, unique ID assigned to that object, enabling the system to maintain its identity over time. By incorporating appearance information, DeepSORT provides a more sophisticated approach to data association compared to SORT, which primarily relies on motion-based predictions.

\subsection{Operational Principles}

The operational efficacy of DeepSORT hinges on the synergistic utilization of robust mathematical tools, notably the Kalman filter for state estimation and prediction, and a sophisticated cost function that combines motion-based and appearance-based metrics for accurate data association across frames. Within the DeepSORT algorithm, each tracked object is modeled by a dynamic state vector that encapsulates its kinematic properties. This state vector typically includes information about the object's position, size, aspect ratio, and their respective velocities, allowing the system to predict the object's future location and characteristics.

\subsubsection{Kalman Filter for State Estimation}

The Kalman filter, a recursive probabilistic estimator, plays a pivotal role in DeepSORT by providing a principled approach to predict the future state of each tracked object and to update this prediction based on new observations (detections) in subsequent frames. The state of an object at time $t$ is represented by the state vector $x_t$, defined as:
$$
x_t = [u, v, s, r, \dot{u}, \dot{v}, \dot{s}]^T
$$
where:
\begin{itemize}
    \item $u, v$: Denote the coordinates of the center of the bounding box in the image plane.
    \item $s$: Represents the area of the bounding box.
    \item $r$: Indicates the aspect ratio (width/height) of the bounding box.
    \item $\dot{u}, \dot{v}, \dot{s}$: Correspond to the respective velocities of the center coordinates and the area of the bounding box. The aspect ratio is often assumed to have a constant velocity of zero in the standard DeepSORT implementation.
\end{itemize}

The Kalman filter operates in two distinct phases: prediction and update. In the prediction step, the filter forecasts the state $x_{t+1}$ of each tracked object at the next time step based on its current state $x_t$ and a predefined motion model:
$$
\hat{x}_{t+1} = F \cdot x_t + w
$$
where $F$ is the state transition matrix that models the object's motion over time, and $w$ represents the process noise, typically modeled as a Gaussian distribution with a specific covariance matrix, accounting for uncertainties in the motion model.

Subsequently, when new detection information (bounding boxes from the object detector) becomes available in the current frame, the update step refines the predicted state $\hat{x}_{t+1}$ to obtain a more accurate estimate of the object's true state. This update is performed based on the measurement (the detected bounding box) and its associated uncertainty, using the Kalman gain, which optimally weights the prediction and the measurement.

\subsubsection{Cost Function for Data Association}

A crucial component of DeepSORT is the carefully designed cost function that quantifies the association between the predicted states of the tracked objects and the newly detected bounding boxes in the current frame. This cost function combines two complementary metrics: motion cost and appearance cost, allowing for a more robust and accurate assignment of detections to existing tracks. The overall cost $c(i, j)$ of associating the $i$-th tracked object with the $j$-th detection is given by:
$$
c(i, j) = \alpha \cdot d_{motion}(i, j) + (1 - \alpha) d_{appearance}(i, j)
$$
where:
\begin{itemize}
    \item $d_{motion}(i, j)$: Represents the motion cost, which measures the compatibility between the predicted state of the $i$-th track (from the Kalman filter) and the $j$-th detected bounding box based on their spatial proximity. This is often quantified using the Mahalanobis distance, which considers the uncertainty in the predicted state. A gating mechanism based on the Mahalanobis distance is typically employed to discard unlikely associations based on motion alone.

    \item $d_{appearance}(i, j)$: Represents the appearance cost, which measures the visual similarity between the $i$-th tracked object and the $j$-th detection. This is computed by extracting a feature vector (embedding) for each detected bounding box using a pre-trained deep convolutional neural network. The appearance cost is then calculated as the cosine distance between the feature vector of the tracked object (typically an average of the features from recent detections belonging to that track) and the feature vector of the current detection. This component is particularly effective in resolving identity ambiguities caused by occlusions or similar-looking objects.

    \item $\alpha$: Is a weighting parameter that balances the relative importance of the motion cost and the appearance cost. The optimal value of $\alpha$ is often determined empirically based on the specific characteristics of the tracking scenario and the dataset.
\end{itemize}

\subsubsection{Hungarian Algorithm for Optimal Assignment}

Once the cost matrix, representing the pairwise costs of associating each tracked object with each detected bounding box, is computed, the Hungarian algorithm is employed to solve the linear assignment problem. The Hungarian algorithm finds the optimal one-to-one assignment between the tracked objects and the detections that minimizes the total cost. This ensures that each detection is assigned to at most one tracked object, and each tracked object is matched with at most one detection in the current frame.

\subsubsection{Track Management}

DeepSORT also incorporates a track management system to handle the creation and deletion of tracks. New tracks are initialized for unassigned detections that do not match any existing tracks. Tracks that are not associated with any detections for a certain number of consecutive frames are considered to have left the scene and are subsequently deleted. This mechanism helps to maintain a consistent set of active tracks and prevents the accumulation of stale or erroneous tracks.

By integrating deep learning-based appearance features with Kalman filtering and a sophisticated cost function optimized by the Hungarian algorithm, DeepSORT achieves robust and accurate multi-object tracking, particularly in challenging scenarios where appearance cues are crucial for maintaining object identities.

\section{ActionCLIP}

\subsection{Introduction}

ActionCLIP, introduced by Wang et al. \cite{wang2021actionclip}, represents a novel and powerful deep learning framework designed for video action recognition by synergistically combining the strengths of computer vision and natural language processing. Building upon the foundational principles of the Contrastive Language-Image Pretraining (CLIP) model \cite{radford2021learning}, ActionCLIP extends CLIP's capabilities to the temporal domain of video. It achieves this by integrating deep neural networks capable of encoding both visual information from video frames and textual descriptions of actions. This innovative approach enables the recognition of actions in videos without the necessity of extensive manual video annotation, offering a significant advantage in terms of scalability and adaptability to novel action categories.

The ActionCLIP model takes as input a video, typically segmented into a series of frames or short video clips, alongside natural language descriptions of the actions it is intended to recognize. These textual descriptions can range from simple phrases like "a person walking" to more complex sentences detailing the ongoing activity, such as "a chef preparing a meal in a kitchen." The model's output is the classification of the action occurring within the video or the assignment of the most relevant action label from the provided textual descriptions to the video content. Fundamentally, ActionCLIP leverages the semantic relationship between visual and textual representations to identify the dominant action portrayed in the video based on the provided textual queries.

\subsection{Operational Principles}

The operational framework of ActionCLIP revolves around the principle of learning a shared embedding space where visual features extracted from video and textual features derived from action descriptions are semantically aligned. This alignment is achieved through a contrastive learning objective that encourages visual and textual representations of the same action to be close to each other in the embedding space, while pushing apart representations of different, unrelated actions. The key components of ActionCLIP's operation include video encoding, text encoding, a shared embedding space, and the contrastive learning mechanism.

\subsubsection{Video Encoding}

Each frame or short clip from the input video is processed by a deep neural network, typically a Convolutional Neural Network (CNN) or a Transformer-based network, to extract salient visual features. This video encoder aims to capture the essential visual elements within the video, such as the involved objects, their motion patterns, the interactions between them, and the overall scene context. The output of this encoding process is a high-dimensional feature vector that encapsulates the visual information relevant to action recognition. Different architectures, such as 3D CNNs or combinations of 2D CNNs with temporal aggregation mechanisms, can be employed as the video encoder depending on the specific implementation and the desired trade-off between computational cost and performance.

\subsubsection{Text Encoding}

Concurrently with video processing, ActionCLIP utilizes a language model, often a Transformer-based model like BERT \cite{devlin2018bert} or GPT \cite{radford2019language}, to encode the natural language descriptions of the actions into feature vectors. The text encoder processes the input text and generates a semantic embedding that captures the meaning and context of the action description. Transformer-based models are well-suited for this task due to their ability to model long-range dependencies in text and capture nuanced semantic relationships between words and phrases.

\subsubsection{Shared Embedding Space}

After extracting features from both the video and the text, ActionCLIP projects these features into a common, multi-modal embedding space. The goal of this shared space is to create a representation where semantically similar visual and textual concepts are located close to each other, regardless of their original modality. This allows for direct comparison and matching between video content and action descriptions based on their embedded representations. The projection into the shared embedding space is typically achieved through linear or non-linear layers learned during the training process.

\subsubsection{Contrastive Learning for Alignment}

To effectively align the visual and textual representations in the shared embedding space, ActionCLIP employs a contrastive learning approach. This learning paradigm aims to maximize the similarity between the feature vectors of corresponding video-text pairs (positive pairs, e.g., a video of someone running and the text "a person running") while minimizing the similarity between feature vectors of non-corresponding pairs (negative pairs, e.g., a video of someone running and the text "a person swimming").

The contrastive loss function used in ActionCLIP can be formally expressed as:
$$
\mathcal{L} = \sum_{i} \log \frac{\exp(\text{sim}(f_i, g_i) / \tau)}{\sum_{j} \exp(\text{sim}(f_i, g_j) / \tau)} + \sum_{i} \log \frac{\exp(\text{sim}(g_i, f_i) / \tau)}{\sum_{j} \exp(\text{sim}(g_i, f_j) / \tau)}
$$
where:
\begin{itemize}
    \item $f_i$: represents the feature vector extracted from the $i$-th video.
    \item $g_i$: represents the feature vector extracted from the corresponding textual description of the action in the $i$-th video.
    \item $\text{sim}(a, b)$: is a similarity function, commonly the cosine similarity, defined as:
    $$
    \text{sim}(a, b) = \frac{a \cdot b}{\|a\| \cdot \|b\|}
    $$
    \item $\tau$: is a temperature parameter that controls the sharpness of the probability distribution.
\end{itemize}
The loss function encourages the similarity between matching video and text embeddings to be high, while pushing the similarity between mismatched pairs to be low. The temperature parameter $\tau$ scales the logits before applying the softmax function, influencing the relative importance of hard negative samples during training.

\subsubsection{Zero-Shot Action Recognition}

A significant advantage of ActionCLIP's architecture and the contrastive learning approach is its ability to perform zero-shot action recognition. Once the model has been trained to align visual and textual representations of a wide range of actions, it can recognize actions in new videos even if those specific actions were not present in the training data. This is achieved by leveraging the semantic understanding captured in the shared embedding space. Given a textual description of a novel action (e.g., "mountain climbing"), the model can compute its text embedding and then compare it to the video embeddings of unseen videos. If a video contains the visual cues corresponding to the "mountain climbing" description, its video embedding will be close to the text embedding of "mountain climbing" in the shared space, allowing for accurate recognition without any explicit training on that specific action. This zero-shot capability makes ActionCLIP highly flexible and adaptable to recognizing a broad spectrum of actions beyond its initial training set.

In summary, ActionCLIP's operational principle relies on learning a semantically rich, shared embedding space for videos and action descriptions through contrastive learning. This enables the model to perform both traditional action classification and, more importantly, zero-shot action recognition by measuring the similarity between visual features of videos and textual features of action labels.


\section{SlowFast Networks}

\subsection{Introduction}

The SlowFast network architecture, proposed by Feichtenhofer et al. \cite{feichtenhofer2019slowfast}, represents a significant advancement in deep learning for video action recognition. It addresses the inherent complexity of human actions, which often involve a combination of slow, deliberate movements and fast, transient changes. The core innovation of SlowFast lies in its dual-stream design, comprising a "slow" stream operating at a lower temporal resolution and a "fast" stream operating at a higher temporal resolution. By simultaneously processing video input through these two streams, the model aims to capture both the semantic context and the fine-grained motion cues that are essential for accurate and robust action recognition across a wide range of temporal dynamics. This approach allows the network to effectively disentangle the processing of what is happening (slow stream) from how it is happening (fast stream).

The input to the SlowFast model is a video sequence, which is sampled into frames and then fed into both the slow and fast processing pathways. The slow stream processes a sparse sampling of the input frames, focusing on capturing the overall scene context, the identity of the actors, and the slowly evolving aspects of the action. In contrast, the fast stream processes a denser sampling of the frames, enabling it to capture the rapid temporal changes, such as the quick movements of limbs or objects, which are crucial for distinguishing between different actions. The output of the SlowFast network is a classification label indicating the recognized action within the video, such as "running," "swimming," "jumping," or more complex activities. The model's strength lies in its ability to fuse the complementary information from the slow and fast streams to achieve high accuracy in classifying actions characterized by diverse temporal scales.

\subsection{Operational Principles}

The SlowFast architecture is predicated on the principle of processing visual information at multiple temporal resolutions, inspired by the human visual system's ability to analyze both the static context and the dynamic motion within a scene. This dual-stream approach allows the network to learn rich spatiotemporal features that are more discriminative for action recognition than those learned by single-stream or purely temporal models.

\subsubsection{Slow Stream: Encoding Spatial Semantics and Long-Term Context}

The slow stream is designed to capture the spatial context and the slowly changing aspects of an action. It operates on a temporally downsampled version of the input video, processing frames at a lower frame rate (e.g., every $T_s$ frames). This lower temporal resolution allows the slow stream to focus on the static elements of the scene, the objects present, and the overall configuration of the environment and actors. The architecture of the slow stream typically involves 3D convolutional layers ($Conv3D$) that extend 2D convolutions into the temporal dimension, enabling the extraction of spatiotemporal features. Let the input video be a sequence of frames $V = \{F_1, F_2, ..., F_N\}$. The slow stream processes a subsequence $V_{slow} = \{F_{1}, F_{1+T_s}, F_{1+2T_s}, ...\}$. The feature maps extracted by the slow stream, $f_{slow}$, can be seen as encoding high-level semantic information and the long-term evolution of the scene.

\subsubsection{Fast Stream: Encoding Fine-Grained Motion Dynamics}

The fast stream, in contrast, is designed to capture the rapid temporal changes and motion dynamics that are characteristic of many actions. It operates on a higher temporal resolution, processing frames at a rate $T_f$, where typically $T_f < T_s$ (often $T_f = 1$, processing all or a larger subset of the original frames). The input to the fast stream is a subsequence $V_{fast} = \{F_{1}, F_{1+T_f}, F_{1+2T_f}, ...\}$. Similar to the slow stream, the fast stream also employs $Conv3D$ layers to extract spatiotemporal features, denoted as $f_{fast}$. However, due to its higher temporal resolution, the fast stream is more sensitive to the short-term movements and rapid changes in the video. To enhance its motion sensitivity while maintaining computational efficiency, the fast stream often uses fewer spatial convolutional filters compared to the slow stream.

\subsubsection{Spatiotemporal Convolutional Operations}

Both the slow and fast streams heavily rely on 3D convolutional operations. A $Conv3D$ kernel of size $t \times h \times w$ operates on a volume of $t$ consecutive frames, each of size $h \times w$, to extract features that are sensitive to both spatial and temporal variations. The output feature map at a specific spatial location $(x, y)$ and temporal location $z$ in the $l$-th layer can be expressed as:

$$
f^{(l)}(x, y, z) = \sigma \left( b^{(l)} + \sum_{c=1}^{C_{l-1}} \sum_{dt=0}^{T_k-1} \sum_{dh=0}^{H_k-1} \sum_{dw=0}^{W_k-1} w^{(l)}_{c, dt, dh, dw} \cdot f^{(l-1)}(x+dw, y+dh, z+dt) \right)
$$

where $\sigma$ is an activation function, $b^{(l)}$ is the bias, $C_{l-1}$ is the number of channels in the $(l-1)$-th layer, and $w^{(l)}_{c, dt, dh, dw}$ are the weights of the $Conv3D$ kernel of temporal size $T_k$, height $H_k$, and width $W_k$.

\subsubsection{Lateral Connections and Feature Fusion}

To effectively combine the information from the slow and fast streams, SlowFast employs lateral connections that allow for the transfer of information between the two streams at various stages of the network. A common approach is to inject features from the slow stream into the fast stream. This fusion allows the fast stream to be conditioned by the contextual information captured by the slow stream. The fusion operation can take various forms, such as addition, concatenation, or more complex learned transformations. For instance, a lateral connection might involve transforming the feature map of the slow stream, $f_{slow}$, through a convolutional layer and then adding it to the feature map of the fast stream, $f_{fast}$, after it has been appropriately resized to match the dimensions:

$$
f'_{fast} = f_{fast} + \mathcal{T}(f_{slow})
$$

where $\mathcal{T}(\cdot)$ represents a series of convolutional and potentially upsampling operations to align the spatial and temporal dimensions of $f_{slow}$ with $f_{fast}$.

The final features from the slow and fast streams are then typically concatenated or combined through other fusion mechanisms (e.g., element-wise multiplication followed by convolution) before being passed to the classification head, which usually consists of fully connected layers and a softmax activation function to produce the probability distribution over the action classes:

$$
P(\text{action} = c | f_{slow}, f_{fast}) = \text{softmax}(W \cdot \text{Fusion}(f_{slow}, f_{fast}) + b)
$$

where $\text{Fusion}(\cdot, \cdot)$ represents the feature fusion operation, $W$ and $b$ are the weights and bias of the final classification layer, and $c$ is one of the action classes.

Through the synergistic operation of the slow and fast streams, coupled with effective feature fusion, the SlowFast network learns robust spatiotemporal representations that are highly effective for recognizing a wide range of human actions in video. The dual-stream architecture allows the model to capture both the essential context and the critical motion cues necessary for accurate action classification.

\section{Supervised Contrastive Learning for Enhanced Representation Learning}

\subsection{Introduction}

Supervised Contrastive Learning (SCL), as introduced by Khosla et al. \cite{khosla2020supervised}, is a powerful machine learning technique within the realm of deep learning, specifically designed to enhance the process of representation learning for supervised classification tasks. Building upon the foundations of Contrastive Learning (CL) \cite{hadsell2006dimensionality, chen2020simple}, SCL leverages the availability of class labels to create more discriminative feature embeddings. By explicitly contrasting data points based on their assigned labels, SCL enables models to learn feature spaces where instances of the same class are tightly clustered together, while instances from different classes are well-separated.

The primary objective of Supervised Contrastive Learning is to learn feature representations such that data samples belonging to the same class (sharing the same label) exhibit high similarity and thus are located close to each other in the embedding space. Conversely, data samples from different classes (having different labels) are pushed apart, resulting in low similarity and greater distance in the embedding space. This learning objective leads to the development of a robust and highly discriminative feature space, where the data points of each class form compact clusters that are well-isolated from other classes, ultimately improving the performance of downstream classification tasks.

\subsection{Operational Principles}

To achieve the goal of learning discriminative feature representations, Supervised Contrastive Learning employs a contrastive loss function that is adapted to utilize the available label information. This loss function compares data points within the embedding space and optimizes the embeddings such that points with the same label are drawn closer together, while points with different labels are pushed farther apart. The Supervised Contrastive Loss can be mathematically expressed as follows:

$$
\mathcal{L}_{SCL} = \sum_{i \in N} \frac{-1}{|P(i)|} \sum_{p \in P(i)} \log \frac{\exp(\text{sim}(z_i, z_p) / \tau)}{\sum_{a \in A(i)} \exp(\text{sim}(z_i, z_a) / \tau) + \exp(\text{sim}(z_i, z_i) / \tau)}
$$

where:
\begin{itemize}
    \item $N$ is the total number of samples in the batch.
    \item $P(i)$ is the set of indices of all data points in the batch that have the same label as the $i$-th data point, excluding $i$ itself. Thus, $|P(i)|$ is the number of positive pairs for the $i$-th sample within the batch.
    \item $A(i)$ is the set of indices of all data points in the batch that have a different label from the $i$-th data point. These represent the negative samples for the $i$-th data point within the batch.
    \item $\text{sim}(z_i, z_j)$ is a similarity function that measures the similarity between the feature vectors $z_i$ and $z_j$ of samples $i$ and $j$. Cosine similarity is a commonly used similarity metric:
    $$
    \text{sim}(z_i, z_j) = \frac{z_i \cdot z_j}{\|z_i\| \cdot \|z_j\|}
    $$
    where $z_i$ and $z_j$ are typically L2 normalized.
    \item $\tau$ is the temperature parameter, a positive scalar that controls the sharpness of the contrastive distribution. A lower temperature makes the loss more sensitive to hard negative samples.
    \item $z_i$ is the feature vector (embedding) of the $i$-th data sample, generated by a deep neural network.
\end{itemize}

The Supervised Contrastive Loss operates on the principle of encouraging clustering of samples with the same label while promoting separation between samples with different labels within the embedding space. For each anchor point $i$ in the batch, the loss considers all other points with the same label as positive samples and all points with different labels as negative samples. The objective is to maximize the similarity between the anchor and its positive samples (numerator of the log term) while minimizing the similarity between the anchor and its negative samples (denominator of the log term). The normalization by $|P(i)|$ averages the loss over all positive pairs for each anchor. The inclusion of $\exp(\text{sim}(z_i, z_i) / \tau)$ in the denominator accounts for the similarity of the sample with itself, although its contribution is constant and primarily ensures numerical stability.

\subsection{Comparison with Unsupervised Contrastive Learning}

Supervised Contrastive Learning builds upon the principles of unsupervised Contrastive Learning (UCL) methods like SimCLR \cite{chen2020simple} and MoCo \cite{he2020momentum}, but with a crucial distinction: the utilization of label information. Here's a comparison of the two paradigms:

\subsubsection{Unsupervised Contrastive Learning}

\textbf{Data Augmentation as Supervision:} UCL relies on generating multiple augmented views of the same unlabeled data instance to create positive pairs. The underlying assumption is that different augmented views of the same instance should have similar representations. Negative pairs are formed by augmented views of different data instances within the batch.
\textbf{Loss Function:} The loss function in UCL aims to maximize the similarity between different augmented views of the same instance and minimize the similarity between augmented views of different instances. A common form of the loss for an anchor $i$ and its positive counterpart $j$ is:
    $$
    \mathcal{L}_{UCL}^{(i)} = - \log \frac{\exp(\text{sim}(z_i, z_j) / \tau)}{\sum_{k \neq i} \exp(\text{sim}(z_i, z_k) / \tau)}
    $$
    where $z_k$ represents the embeddings of all other augmented views in the batch (including other positive pairs).
\textbf{Objective:} The primary objective is to learn general-purpose representations that capture the inherent structure and variations within the unlabeled data. These representations can then be used for various downstream tasks, often requiring fine-tuning with labeled data.

\subsubsection{Supervised Contrastive Learning}

\textbf{Label-Based Positive Pairs:} SCL directly utilizes class labels to define positive pairs. Instead of relying solely on augmentations of the same instance, SCL considers all instances within the same batch that share the same class label as positive pairs for a given anchor instance.

\textbf{More Informative Negative Samples:} The negative samples in SCL are explicitly defined as instances belonging to different classes within the batch. This provides more direct and informative negative examples for the model to discriminate against.

\textbf{Loss Function:} As shown in the previous section, the Supervised Contrastive Loss explicitly sums over all positive pairs for each anchor and contrasts them against all negative samples in the batch.

\textbf{Objective:} The primary objective of SCL is to learn representations that are specifically tailored for classification. By leveraging label information during training, SCL aims to create a feature space where class separability is maximized, leading to improved performance on supervised classification tasks, often with less reliance on extensive data augmentation compared to UCL for downstream tasks.

\textbf{Key Differences and Advantages of SCL:}

\textbf{Stronger Supervision:} SCL benefits from the explicit supervision provided by class labels, leading to more discriminative representations compared to UCL, which relies on the inductive bias that different augmentations of the same instance should be similar.
\textbf{Improved Class Separability:} By explicitly pulling together instances of the same class and pushing apart instances of different classes, SCL directly optimizes for class separability in the embedding space.
\textbf{Potentially Less Reliance on Extensive Augmentations:} While augmentations can still be beneficial in SCL to improve robustness, the strong signal from the labels can reduce the need for very diverse and aggressive augmentations that are often crucial for the success of UCL.
\textbf{Task-Specific Representations:} The representations learned by SCL are more directly aligned with the classification task for which the labels are available, often leading to better performance on these tasks compared to the more general-purpose representations learned by UCL.

In conclusion, Supervised Contrastive Learning represents a significant advancement over unsupervised contrastive learning for supervised tasks by effectively incorporating label information into the contrastive learning framework. This leads to the learning of more discriminative and task-specific feature representations, ultimately enhancing the performance of classification models.

